\begin{abstract}
Byzantine-robust aggregation inspects individual client updates, while homomorphic encryption exists to hide them. We measure two orderings over identical updates in a single-process simulator with real CKKS arithmetic: Encrypt-then-Filter (EtF), where a key holder sees pairwise distances, and Filter-then-Encrypt (FTE), where a filter sees plaintext updates. On every supported record, both selected the same clients as the plaintext rule. EtF Krum took about 25 times as long per round as FTE Krum on FEMNIST and about 10 times as long on Speech Commands, but FTE's filter sees every update, and only FTE can run trimmed mean and Bulyan. In 120 FEMNIST runs (three seeds, eight conditions, five rules) no rule was robust overall. Robust rules lost 0.03 to 0.28 accuracy without an attack, only scaling gave any of them an advantage over FedAvg, and a zero-update attack cut Krum from 0.478 to 0.054. Multi-Krum led the robust rules on accuracy but gave no protection against a DBA-style backdoor. On Speech Commands (BC-ResNet-8, three seeds, 40 rounds, 200 speakers, not converged) 10 of 36 runs diverged, all under FedAvg or Multi-Krum, and trimmed mean stopped a feature-space backdoor on only one seed of three. We conclude that the ordering decides cost and exposure, and the rule's fit to the threat decides robustness. FTE's plaintext filter is a real trust cost, not a free speed-up.
\end{abstract}
